\section{An exact weight-eight commutator normalization}
\label{sec:weight-eight-normalization}

Put $M=\ker\lambda$, and write
\[
 \begin{gathered}
 D=N_{\mathrm{rad}},\qquad D_3=D\cap\ker m_3,
 \qquad D_5=D_3\cap\ker m_5,\\
 D_7=D_5\cap\ker m_7,\qquad J_7=D_7\cap H.
\end{gathered}
\]
The preceding normalizations give $G=D_7H$. All Magnus expansions
below are in the free pro-$2$ group, with $x=1+S$ and $y=1+T$.
For $g=(1,f)\in D_7$, define the mixed coefficient
\[
 m_8(g)=[S^4TST^2]f.
\]
The depth-one coefficient $[S^7T]f$ is zero; the first available
degree-eight direction is an Ihara bracket.

\begin{theorem}\label{thm:weight-eight-exact-image}
Every word belonging to $D_7$ starts in degree eight. Its homogeneous
part of degree eight lies in a primitive rank-one integral Lie module,
and
\[
 m_8(D_7)=m_8(J_7)=\mathbf Z_2.
\]
Set $D_8=D_7\cap\ker m_8$ and $J_8=D_8\cap H$. Then $D_8$ is a
proper closed normal subgroup of $G$,
\[
 D_7=D_8J_7,\qquad G=D_8H,\qquad
 B\simeq D_8/\overline{[G,D_8]J_8}.
\]
Every word belonging to $D_8$ starts in degree at least nine.
The continuous character $\eta_8=m_8\bmod2$ is $G$-invariant on
$D_7$ and has nonzero arithmetic restriction. It has no
$G$-invariant extension to $D$ or to $N=\ker(\lambda,\rho_2)$.
\end{theorem}

We give the integral leading calculation and an arithmetic
commutator. No existence claim is drawn from a truncated system of
necessary equations.

\subsection{The integral leading line}

For a Lyndon word $w$ in $S<T$, let $L_w$ be its standard Lie
bracket, using the longest proper Lyndon suffix. Write
\[
 \psi_3=L_{SST}-L_{STT},
\]
and, in the degree-five Lyndon order, put
\[
 \psi_5=L_{SSSST}-2L_{SSSTT}
       +\tfrac32 L_{SSTST}+2L_{SSTTT}
       +\tfrac12 L_{STSTT}-L_{STTTT}.
\]
For a Lie element $P$, define the derivation
$\partial_P(S)=0$, $\partial_P(T)=[T,P]$, and use the Ihara bracket
\[
 \{P,Q\}=[P,Q]+\partial_P(Q)-\partial_Q(P).
\]
Let
\[
 \mathcal L_8=2\{\psi_3,\psi_5\}.
\]
Its integral Lyndon coordinates are as follows:
\[
\begin{array}{cr@{\qquad}cr@{\qquad}cr}
 w&[L_w]\mathcal L_8&w&[L_w]\mathcal L_8&w&[L_w]\mathcal L_8\\ \hline
 SSSSSSST&0&SSSTSTTT&-6&SSTTSTST&21\\
 SSSSSSTT&0&SSSTTSST&4&SSTTSTTT&3\\
 SSSSSTST&-4&SSSTTSTT&-3&SSTTTSTT&-15\\
 SSSSSTTT&0&SSSTTTST&-6&SSTTTTST&-3\\
 SSSSTSST&-2&SSSTTTTT&0&SSTTTTTT&0\\
 SSSSTSTT&7&SSTSSTST&-5&STSTSTTT&8\\
 SSSSTTST&17&SSTSSTTT&3&STSTTSTT&-11\\
 SSSSTTTT&0&SSTSTSTT&3&STSTTTTT&-4\\
 SSSTSSTT&-1&SSTSTTST&6&STTSTTTT&-10\\
 SSSTSTST&-12&SSTSTTTT&7&STTTTTTT&0
\end{array}
\]
In particular $\mathcal L_8$ is primitive. Direct associative
expansion gives
\begin{equation}\label{eq:weight-eight-mixed-coefficient}
 [S^4TST^2]\mathcal L_8=7,
 \qquad [S^7T]\mathcal L_8=0.
\end{equation}
The bracket formula and the displayed finite table specify this
identity without a choice of a group lift of either rational Lie
normalization.

\begin{lemma}\label{lem:weight-eight-leading-line}
Suppose $f\in\gamma_8$ satisfies the two-cycle and three-cycle
relations
\[
 f(y,x)f(x,y)=1,\qquad
 f(z,x)f(y,z)f(x,y)=1,\qquad z=(xy)^{-1}.
\]
Then its homogeneous part of degree eight is
\[
 \operatorname{in}_8(f)=b\mathcal L_8\quad\text{for a unique }b\in\mathbf Z_2.
\]
If $h_w$ denotes a degree-eight Lyndon coordinate, then
$b=-h_{SSSTSSTT}$ and $m_8(f)=7b$.
\end{lemma}

\begin{proof}
The following degree-nine certificate proves the necessary rational
line. Its coordinate $-1$ then proves the stated integral assertion.
Collect $f$ modulo $\gamma_{10}$ in the standard Lyndon group
commutators $C_w$, with integral coordinates $h_w$. Explicitly, for the
standard factorization $w=uv$, put
\[
 C_S=x,\qquad C_T=y,\qquad C_w=C_uC_vC_u^{-1}C_v^{-1}.
\]
Set
\[
 F=1+\sum_{8\le |w|\le9}h_w(C_w-1).
\]
There are $30+56=86$ columns, in order of weight and then lexicographic
order. The expression for $F$ is exact through degree nine, since
products of two nonconstant factors have degree at least sixteen.
All the following polynomial operations are truncated after degree
nine. Put $X=1+S$, $Y=1+T$, $Z=Y^{-1}X^{-1}$, and form
\[
 I_v=[v]\bigl(F(Y,X)+F(X,Y)-2\bigr),\qquad
 II_v=[v]\bigl(F(Z,X)+F(Y,Z)+F(X,Y)-3\bigr).
\]
These are integral linear forms on the $86$ Hall coordinates and
vanish for the given word.

Encode $I_v$ by $2^{|v|}+\operatorname{bin}(v)$ and $II_v$ by
$1024+2^{|v|}+\operatorname{bin}(v)$, with $S=0$, $T=1$. Let $A_R$
be the $66$ by $86$ matrix consisting of the following rows, in the
written order:
\[
\begin{aligned}
R=( &257,259,261,263,265,267,269,275,277,281,283,285,293,301,\\
    &513,515,517,519,521,523,525,527,529,531,533,535,537,539,\\
    &541,547,549,551,553,555,557,563,565,569,587,589,595,597,\\
    &1281,1287,1291,1299,1301,1317,1537,1539,1541,1545,\\
    &1547,1549,1553,1555,1557,1559,1561,1563,1571,1573,\\
    &1577,1579,1581,1589).
\end{aligned}
\]
Let $C$ be its first $66$ columns and let $A_{R,C}$ denote the
corresponding square matrix. If $\ell_j$ is the $j$th entry of the
thirty-coordinate table above, ordered lexicographically by its word,
define $T$ to have the $29$ rows
\[
 T_j=e_j+\ell_j e_9\quad(1\le j\le30,\ j\ne9),
\]
where the $e_j$ are rows of length $86$. The finite identities are
\begin{equation}\label{eq:weight-eight-leading-certificate}
 \det A_{R,C}=-12312=-2^3\,3^4\,19,
 \qquad
 T_C A_{R,C}^{-1}A_R=T.
\end{equation}
Here every matrix entry is determined by the displayed group
commutator recursion and the finite geometric expansion
$(1+A)^{-1}=\sum_{j=0}^{9}(-A)^j$. Thus
\eqref{eq:weight-eight-leading-certificate} is an explicit finite
rational matrix identity, not a rank assumption. Multiplying
$A_Rh=0$ by $T_C A_{R,C}^{-1}$ gives
$h_j=-\ell_jh_9$ for all $j\le30$. This proves the claimed line
over $\mathbf Q_2$. Since $h_9\in\mathbf Z_2$ and $\ell_9=-1$,
its scalar $b=-h_9$ is integral. The mixed coefficient is then
$7b$ by \eqref{eq:weight-eight-mixed-coefficient}.
\end{proof}

The earlier leading degree-seven calculation shows that a word in
$D_5$ has its degree-seven part determined by $m_7$. Therefore every
word in $D_7$ starts in degree eight and satisfies the preceding
lemma. In particular its degree-eight part vanishes exactly when
$m_8$ vanishes. This upper-lattice argument uses only the two-cycle
and three-cycle identities once the preceding normalization has
been made; it uses no new pentagon coefficient.

\subsection{A primitive arithmetic commutator}

Recall the three homomorphisms on all of $M$, together with their
scalar conjugation laws:
\[
 \begin{array}{c|c|c}
  \text{homomorphism}&\text{formula}&\text{conjugation weight}\\ \hline
  m_3 &[S^2T]f&\lambda_2^3\\
  \tau_5&12[S^4T]f-11m_3&\lambda_2^5\\
  \tau_7&360[S^6T]f+900[S^5T]f+780[S^4T]f-239m_3
            &\lambda_2^7.
 \end{array}
\]
They are the coefficients of the logarithmic Alexander beta
cocycle: $m_3=c_3$, $\tau_5=12c_5$, and $\tau_7=360c_7$, where
$c_j=[u^{j-1}v]\log B(e^u,e^v)$. Thus these formulas apply before
any radical normalization. On $D_3$ one has $\tau_5=12m_5$, and
on $D_5$ one has $\tau_7=360m_7$.

\begin{lemma}\label{lem:weight-eight-arithmetic-commutator}
One has $\overline{[M,M]}\subseteq D_7$. There are
$\sigma\in H\cap M$ and $h\in H$ such that, with
$s=h\sigma h^{-1}$,
\[
 w_8=[\sigma,s]\in J_7,
 \qquad m_8(w_8)=63\,m_3(\sigma)\,\tau_5(\sigma)
                 \in\mathbf Z_2^\times.
\]
Moreover $w_8=[\sigma,k]$ for some $k\in D\cap H$.
\end{lemma}

\begin{proof}
Every rational-prime Kummer cocycle is a homomorphism on $M$.
Consequently $[M,M]\subseteq D$. The homomorphism $m_3$ makes
these commutators belong to $D_3$; the homomorphism $\tau_5$, which
there equals $12m_5$, makes them belong to $D_5$; and $\tau_7=360m_7$
on $D_5$ makes them belong to $D_7$. All the kernels are closed,
which proves the first assertion.

The exact arithmetic Kummer image supplies
$\sigma\in G_{\mathbf Q^{\mathrm{ab}}}\subset H\cap M$ with
\[
 \rho_2(\sigma)=2\in\widehat{\mathbf Z},\qquad
 \rho_q(\sigma)=0\quad(q\text{ odd prime}).
\]
Write $a=m_3(\sigma)$. The established congruence
$\rho_2\equiv-2m_3\pmod4$ makes $a$ odd, and the formula for
$\tau_5$ makes $\tau_5(\sigma)$ odd as well.
Choose $h\in H$ with $\lambda_2(h)=3$, and put $t=\lambda(h)$.
Profinite powers are defined by the continuous map
$\widehat{\mathbf Z}\to\overline{\langle\sigma\rangle}$ sending
$1$ to $\sigma$. Hence both
\[
 k=s\sigma^{-t}\quad\text{and}\quad r=s\sigma^{-27}
\]
are elements of $H\cap M$. Conjugation has weight $\lambda$
on every Kummer homomorphism, so $k\in D\cap H$. Since powers of
$\sigma$ commute with $\sigma$,
\[
 [\sigma,k]=[\sigma,s]=[\sigma,r]=w_8.
\]
The weight-three law gives $m_3(r)=(3^3-27)a=0$. The low-degree
cycle calculation, whose degree-three and degree-four Hall
coordinates are $(m_3,-m_3)$ and $(-m_3,m_3,-m_3)$, therefore shows
that $r$ starts in degree five. Its leading part is a multiple of
$\psi_5$, by the preceding integral degree-five calculation. The
weight-five law gives its multiplier explicitly:
\[
 [S^4T]f_r=\frac{3^5-27}{12}\tau_5(\sigma)
       =18\tau_5(\sigma).
\]
The leading part of $\sigma$ is $a\psi_3$.

For completeness, with the normalized action
$\Phi_f(x)=x$, $\Phi_f(y)=f^{-1}yf$, multiplication of parameters
is $f\star g=f\Phi_f(g)$. Its leading Lie bracket is
$[P,Q]+\partial_P(Q)-\partial_Q(P)$, since the first-order change
in $T$ is $[T,P]$. Thus the leading part of this group
commutator is
\[
 \operatorname{in}_8(w_8)
   =18a\tau_5(\sigma)\{\psi_3,\psi_5\}
   =9a\tau_5(\sigma)\mathcal L_8.
\]
Taking the coefficient in
\eqref{eq:weight-eight-mixed-coefficient} gives the claimed odd
value $63a\tau_5(\sigma)$.
\end{proof}

\begin{proof}[Proof of Theorem~\ref{thm:weight-eight-exact-image}]
The leading coefficient $m_8$ is an additive continuous map on
$D_7$. This follows either by multiplying the filtered normalized
automorphisms or by noting that substitutions affect a leading
word only above degree eight. Conjugation acts on it by
\[
 m_8(gng^{-1})=\lambda_2(g)^8m_8(n)
 \quad(g\in G,\ n\in D_7).
\]
Indeed the associated graded action on the two free generators is
scalar $\lambda_2$, and conjugation of a tangential derivation of
degree eight has weight $\lambda_2^8$.

The image of $J_7$ is a compact additive subgroup of $\mathbf Z_2$
containing the unit given by the arithmetic commutator. It is
therefore all of $\mathbf Z_2$; the same holds for $D_7$.
The scalar law proves that $D_8$ is normal and that $\eta_8$ is
invariant. The unit proves properness and arithmetic nonvanishing.
Equality of the arithmetic and full images gives $D_7=D_8J_7$, hence
$G=D_8H$. The leading-line lemma shows that words in $D_8$ start
in degree at least nine.

The relative abelian quotient is unchanged by this replacement.
Explicitly, a continuous $G$-invariant homomorphism from $D_8$ to
a profinite abelian group, trivial on $J_8$, extends uniquely to
$G=D_8H$ by sending $dh$ to the value at $d$. The intersection
condition proves well-definedness, and the compact surjection
$D_8\times H\to G$ proves continuity. This is exactly the universal
property of $D_8/\overline{[G,D_8]J_8}$ and of the preceding quotient
$B$.

Finally let $\chi:D\to\mathbf F_2$ be $G$-invariant. The arithmetic
word satisfies $w_8=[\sigma,k]$ with $k\in D$, so
\[
 \chi(w_8)=\chi(\sigma k\sigma^{-1})-\chi(k)=0,
 \qquad \eta_8(w_8)=1.
\]
Thus $\chi$ cannot extend $\eta_8$. An extension to $N$ would
restrict to one on $D$, so it is impossible as well.
\end{proof}

\begin{corollary}\label{cor:weight-eight-original-character-kernel}
Let $E$ be any closed $G$-stable subgroup containing $D$ and contained
in $M$, and let $A$ be a finite abelian group with trivial $G$-action.
Then
\[
 \{\chi\in\operatorname{Hom}_{\mathrm{cts}}(E,A)^G:
                         \chi(D_8)=0\}
 =
 \{\chi\in\operatorname{Hom}_{\mathrm{cts}}(E,A)^G:
                         \chi(D_7)=0\}.
\]
In particular,
\[
 \ker\!\left(H^1(N,\mathbf F_2)^G
                 \longrightarrow H^1(D_8,\mathbf F_2)\right)
 =\left(\bigoplus_{q\text{ odd prime}}\mathbf F_2\psi_q\right)
       \oplus\mathbf F_2\eta_3.
\]
\end{corollary}

\begin{proof}
Every invariant homomorphism on $E$ kills $w_8=[\sigma,k]$, since
$k\in D\subseteq E$. If it also kills $D_8$, its restriction to
$D_7$ factors through $D_7/D_8\simeq\mathbf Z_2$. The image of
$w_8$ topologically generates this quotient, so the restriction is
zero. The reverse inclusion is immediate. Apply this assertion
to $E=N$, $A=\mathbf F_2$, and use the previously proved kernel
formula for restriction to $D_7$.
\end{proof}

\subsection{The extension threshold of the eighth character}
\label{subsec:eta-eight-extension-threshold}

Complex conjugation gives an arithmetic obstruction to extending the
new character through the preceding normalizations.

\begin{proposition}\label{prop:eta-eight-no-extension-d3}
The character $\eta_8$ on $D_7$ extends to the invariant character
\[
 \alpha_8:D_5\longrightarrow\mathbf F_2,\qquad
 \alpha_8(1,f)=[S^4TST^2]f\pmod2,
\]
but has no $G$-invariant extension to $D_3$. More precisely, no
nonzero linear combination of $\eta_7$ and $\alpha_8$ on $D_5$
extends to a $G$-invariant character of $D_3$.
\end{proposition}


\begin{proof}
The exact weight-seven theorem puts the degree-seven part of every
$D_5$ word in $8\mathbf Z_2\mathcal L_7$, where $\mathcal L_7$
is the primitive vector of \eqref{eq:weight-seven-primitive-vector}.
Thus its Magnus expansion over $\mathbf F_2$ starts in degree eight.
The first nonzero term is additive under the normalized GT product,
and conjugation acts by $\lambda_2^8$, which is one modulo two.
This follows by substituting the linear terms in the GT automorphism:
higher terms raise degree, and a homogeneous ambiguity commuting with
$T$ would be a pure power of $T$, absent from a commutator word.
Hence $\alpha_8$ is an invariant character and restricts to $\eta_8$.

On $M=\ker\lambda$ use the additive coordinates
\[
 a=m_3,\qquad T_5=12[S^4T]f-11a,\qquad
 T_7=360[S^6T]f-1020[S^4T]f+661a.
\]
Their conjugation weights are $\lambda_2^3,\lambda_2^5,\lambda_2^7$.
Together with the rational Kummer homomorphisms they imply
\[
 [M,M]\subset D_7,\qquad [M,D_7]\subset D_8.
\]
For the first inclusion all four families vanish on commutators;
then $T_5=12m_5$ and $T_7=360m_7$, and $\mathbf Z_2$ has no torsion.
For the second inclusion use the weight-eight law on $D_7$ and
$\lambda(M)=1$. Therefore $M/D_8$ has nilpotency class at most two,
and its commutator subgroup is contained in $D_7/D_8\simeq\mathbf Z_2$.
The previously proved commutator formula is
\begin{equation}\label{eq:eta-eight-commutator-pairing}
 m_8([z,z'])=\frac7{24}\bigl(a(z)T_5(z')-a(z')T_5(z)\bigr)
 \qquad(z,z'\in M).
\end{equation}

Let $c\in H$ be complex conjugation, and choose $h\in H$ with
$\lambda_2(h)=3$. Put $t=\lambda(h)$ and $\sigma=[h,c]$.
Since $c^2=1$, $\lambda(c)=-1$ and $\rho_2(c)=0$, one has
\[
 c\sigma c^{-1}=\sigma^{-1},\qquad
 \rho_2(\sigma)=2\rho_2(h).
\]
The established parity law makes $\rho_2(h)$ odd at two, and
$\rho_2(\sigma)\equiv-2m_3(\sigma)\pmod4$ makes
\[
 A=a(\sigma),\qquad B=T_5(\sigma),\qquad C=T_7(\sigma)
\]
all odd. The last two assertions also follow directly from their
formulas above. The odd-prime Kummer coordinates of $\sigma$ need
not vanish.

For $e\in\widehat{\mathbf Z}$ put
$D_e(z)=hzh^{-1}z^{-e}$, and define
\[
 z_1=D_t(\sigma),\qquad z_3=D_{t^3}(z_1),\qquad
 z_5=D_{t^5}(z_3).
\]
The weight laws give $z_1\in D\cap H$, $z_3\in D_3\cap H$ and
$z_5\in D_5\cap H$. Moreover
\[
 m_7(z_5)=\frac{(3^7-3)(3^7-3^3)(3^7-3^5)}{360}C
        =128\cdot199017\,C,
\]
so $\eta_7(z_5)=1$. Every invariant quadratic character on $D_3$
kills $z_5$, since conjugation fixes its value on $z_3$ and
$t^5$ is odd at two.

A second arithmetic element is $v=cz_3c^{-1}z_3$.
Every rational Kummer coordinate and each of $a,T_5,T_7$ changes
sign under $c$, so $v\in D_7$. Every invariant character on $D_3$
also kills $v$. We compute $m_8(v)$ in the class-two quotient.
Put $X=\sigma$, $Y=h\sigma h^{-1}$, $Z=h^2\sigma h^{-2}$.
The identity $chc^{-1}=X^{-1}h$ gives
\[
 cXc^{-1}=X^{-1},\quad
 cYc^{-1}=X^{-1}Y^{-1}X,\quad
 cZc^{-1}=X^{-1}Y^{-1}Z^{-1}YX.
\]
For profinite exponents $p,r$, set $W=ZY^{-p}(YX^{-p})^{-r}$.
Class-two collection gives
\begin{equation}\label{eq:eta-eight-conjugation-norm}
 (cWc^{-1})W=[X,Y]^{pr(p-1)-p-r}
                  [X,Z]^{1-pr}[Y,Z]^{1+p+r}.
\end{equation}
Here is a direct verification. Encode the normal form
$X^aY^bZ^c[X,Y]^d[X,Z]^e[Y,Z]^f$ by $(a,b,c,d,e,f)$.
Multiplication is
\[
 (a,b,c,d,e,f)(A,B,C,D,E,F)
 =(a+A,b+B,c+C,d+D-Ab,e+E-Ac,f+F-Bc),
\]
and its $n$th power has coordinates
\[
 (na,nb,nc,nd-\tbinom n2ab,
             ne-\tbinom n2ac,nf-\tbinom n2bc).
\]
The normal form of $W$ is
\[
 (pr,-p-r,1,\ p^2r+pr(r-1)/2,-pr,p+r).
\]
Conjugation by $c$ sends a coordinate tuple to
$(-a,-b,-c,d+b,e+c,f+c)$. Multiplying these two tuples gives
\eqref{eq:eta-eight-conjugation-norm}. Integer exponent identities
extend to profinite exponents by continuity, including the integral
binomial polynomial $n(n-1)/2$.

For $W=z_3$, the exponents are $p=t,r=t^3$. Only after applying
$m_8$ to the central commutator group do their two-adic values become
$3,27$. Thus the exponents in \eqref{eq:eta-eight-conjugation-norm}
are $132,-80,31$ in that coordinate. Formula
\eqref{eq:eta-eight-commutator-pairing} gives
\[
 m_8([X,Y])=63AB,\quad m_8([X,Z])=17010AB,\quad
 m_8([Y,Z])=413343AB.
\]
Consequently
\[
 m_8(v)=(132\cdot63-80\cdot17010+31\cdot413343)AB
       =11461149\,AB
\]
is odd. Hence $(\eta_7(v),\alpha_8(v))=(0,1)$.

Suppose an invariant character on $D_3$ restricts to
$\mu\eta_7+\nu\alpha_8$ on $D_5$. Evaluation at $v$ forces
$\nu=0$; evaluation at $z_5$ then forces $\mu=0$.
The value $\alpha_8(z_5)$ is unnecessary. Finally, an extension
of $\eta_8$ to $D_3$ would restrict on $D_5$ to
$\alpha_8+\epsilon\eta_7$, because $D_5/D_7\simeq128\mathbf Z_2$.
Both possibilities have been excluded.
\end{proof}

\begin{corollary}\label{cor:d3-character-kernel-d8}
One has
\[
 \ker\!\left(H^1(D_3,\mathbf F_2)^G
                    \longrightarrow H^1(D_8,\mathbf F_2)\right)
       =\mathbf F_2\eta_5.
\]
\end{corollary}
\begin{proof}
Let $\chi$ be an invariant character on $D_3$ that vanishes on $D_8$.
Since $D_7/D_8\simeq\mathbf Z_2$ through $m_8$, its restriction
to $D_7$ is zero or $\eta_8$. Subtracting the corresponding multiple
of $\alpha_8$ from $\chi|_{D_5}$ makes it vanish on $D_7$.
It then factors through $D_5/D_7\simeq128\mathbf Z_2$, so
$\chi|_{D_5}=\mu\eta_7+\nu\alpha_8$.
The proposition forces both coefficients to be zero.
Thus $\chi$ factors through $D_3/D_5\simeq32\mathbf Z_2$ and is
a multiple of $\eta_5$. Conversely $\eta_5$ is invariant on $D_3$
and vanishes on $D_5$, hence on $D_8$.
\end{proof}

